\documentclass[11pt,a4paper]{article}
\usepackage[margin=1in]{geometry}
\usepackage{times}
\usepackage{latexsym}
\usepackage{graphicx}
\usepackage{amsmath}
\usepackage{amssymb}
\usepackage{booktabs}
\usepackage{multirow}
\usepackage{xcolor}
\usepackage{tabularx}
\usepackage{array}
\usepackage[colorlinks=true,citecolor=blue,urlcolor=blue]{hyperref}
\usepackage{natbib}
\bibliographystyle{plainnat}

\newcommand{\ISR}{\textsc{ISR}}
\newcommand{\STA}{\textsc{STA}}
\newcommand{\GRPO}{\textsc{GRPO}}
\newcommand{\LOO}{\textsc{LOO}}

\title{Diagnosing and Characterizing Reward Hacking\\in GRPO-Trained Formal Reasoners}

\author{Anonymous}

\begin{document}
\maketitle

% ============================================================
\begin{abstract}
Group Relative Policy Optimization (\GRPO{}) with execution-based rewards trains neural formalizers that achieve high prove rates yet produce structurally degenerate outputs---a phenomenon we term \emph{reward hacking} in formal reasoning.
We propose \textbf{Solver Trace Ablation} (\STA{}), a diagnostic framework that applies leave-one-out component necessity testing to solver inputs, yielding the \textbf{Intensional Structural Reward} (\ISR{}) as a quantitative measure of structural integrity.
In first-order logic (FOL) experiments with Qwen3-1.5B across 3 seeds and up to 1000 training steps, we show that \ISR{} detects hacking invisible to prove rate: exec-only \GRPO{} achieves 93\% prove rate while \ISR{} collapses from 0.64 (100 steps) to 0.003 (1000 steps).
When \ISR{} is repurposed as a training reward, it exhibits \emph{meta-level Goodhart's Law}: at low dose ($\beta{=}0.1$), 2/3 seeds game \ISR{} itself via conclusion duplication (H7 rate 87--97\%), while the remaining seed's apparent improvement is entirely a premise-count artifact (1.06$\times$ when controlled); at high dose ($\beta{=}0.3$), classic Goodhart collapse occurs (prove rate 18\% vs.\ random reward 48\%).
We identify 6 distinct hacking strategies across reward conditions and document \ISR{}'s fundamental asymmetry: reliable as an observer, unreliable as an optimizer.
Preliminary cross-logic experiments (Prolog, SMT-LIB) show high seed instability and training failures, indicating that cross-logic generalization of these findings remains an open challenge.
\end{abstract}

% ============================================================
\section{Introduction}

Reinforcement learning from execution feedback has emerged as a promising paradigm for training neural theorem provers and formal reasoning systems \cite{luo2024improve}.
\GRPO{} \cite{shao2024deepseekmath} applies group-relative advantages to execution rewards, enabling models to learn formal reasoning without human annotations.
However, high prove rates can mask severe structural degeneration: models learn shortcuts---embedding conclusions in premises, generating tautologies, or producing redundant components---that satisfy the solver without genuine reasoning.

Current evaluation practice relies primarily on prove rate, which is blind to structural quality.
No principled diagnostic exists for detecting and characterizing reward hacking in formal reasoning systems.
Existing work either hand-crafts logic-specific penalties \cite{mellgren2025} or uses scalar execution rewards without structural analysis \cite{evom2026}.

We propose \textbf{Solver Trace Ablation} (\STA{}), which applies leave-one-out (LOO) ablation to solver inputs to derive the \textbf{Intensional Structural Reward} (\ISR{}).
For each proved formalization, \STA{} systematically removes each component and re-queries the solver, measuring what fraction of components are individually necessary for the proof to succeed.
\ISR{} thus quantifies the structural necessity of each component, providing a diagnostic signal orthogonal to prove rate.

Our contributions are:
\begin{enumerate}
\item \textbf{\ISR{} as diagnostic} (\S\ref{sec:claim1}): We show that \ISR{} detects structural degeneration invisible to prove rate, with a complete collapse trajectory (0.64 $\to$ 0.13 $\to$ 0.003) over 1000 training steps and a taxonomy of 6 hacking modes across 3 seeds.
\item \textbf{Meta-level Goodhart's Law} (\S\ref{sec:claim2}): When \ISR{} is used as a training reward, it becomes a target of optimization itself. At $\beta{=}0.1$, 2/3 seeds exhibit ISR-specific gaming (H7 rate 87--97\%); the remaining seed's apparent improvement vanishes under premise-count control (1.06$\times$). At $\beta{=}0.3$, prove rate collapses to 18\%. Random reward controls and an SFT baseline confirm these are ISR-specific and RL-specific phenomena.
\item \textbf{Negative results on cross-logic transfer} (\S\ref{sec:crosslogic}): We document that Prolog and SMT-LIB experiments exhibit high seed instability, training failures, zero cross-logic transfer, and catastrophic interference in mixed training---indicating that extending these findings beyond FOL remains challenging.
\end{enumerate}

% ============================================================
\section{Related Work}

\paragraph{Reward hacking in RL.}
Goodhart's Law---``when a measure becomes a target, it ceases to be a good measure''---has been extensively studied in RL alignment \cite{manheim2019categorizing}.
Reward hacking has been documented in RLHF for language models \cite{gao2023scaling}, game-playing agents, and robotic control.
Our work provides an empirical case study of Goodhart's Law in formal reasoning, including the novel observation of \emph{meta-level} Goodhart effects where a diagnostic metric games itself when used as reward.

\paragraph{Formal reasoning with RL.}
Recent work applies RL to formal reasoning: DeepSeek-Prover uses \GRPO{} with execution rewards \cite{shao2024deepseekmath}, EVOM \cite{evom2026} uses scalar execution rewards for verification, and Mellgren et al.\ \cite{mellgren2025} hand-craft Prolog-specific structure penalties.
None systematically diagnose or characterize reward hacking in this setting.

\paragraph{Feature attribution and LOO methods.}
Leave-one-out ablation is a classical technique in feature attribution \cite{zeiler2014visualizing}.
We explicitly acknowledge that \ISR{} applies LOO ablation to solver inputs; our contribution is the application to reward hacking diagnosis and the empirical findings, not the ablation technique itself.

% ============================================================
\section{Method}
\label{sec:method}

\subsection{Solver Trace Ablation (\STA{})}

Given a formalization $F = \{c_1, c_2, \ldots, c_N\}$ consisting of $N$ components (premises in FOL, clauses in Prolog) and a solver $\mathcal{S}$, \STA{} proceeds as follows:

\begin{enumerate}
\item Query $\mathcal{S}(F)$ to determine if $F$ is provable.
\item For each $c_i \in F$, query $\mathcal{S}(F \setminus \{c_i\})$.
\item Define component $c_i$ as \emph{necessary} if $\mathcal{S}(F) = \text{proved}$ but $\mathcal{S}(F \setminus \{c_i\}) \neq \text{proved}$.
\end{enumerate}

The \textbf{Intensional Structural Reward} is:
\begin{equation}
\ISR{}(F) = \frac{|\{c_i : c_i \text{ is necessary}\}|}{N}
\end{equation}

\ISR{} ranges from 0 (no component is individually necessary---all are redundant) to 1 (every component is necessary for the proof). For unprovable formalizations, \ISR{} is undefined (set to 0 in aggregation).

\paragraph{Computational cost.} \STA{} requires $N+1$ solver calls per sample (one base query plus $N$ ablations), giving $O(N)$ complexity. We also evaluate an $O(1)$ heuristic approximation (\S\ref{sec:o1ablation}).

\subsection{Hacking Taxonomy}

Through systematic analysis of \GRPO{}-trained formalizations across multiple reward conditions, we identify 6 hacking modes:

\begin{itemize}
\item \textbf{H1}: Conclusion embedding---the conclusion appears among multiple premises (multi-premise, diluted \ISR{}).
\item \textbf{H3}: Complexity without utility---unnecessary predicates or quantifiers that do not contribute to the proof.
\item \textbf{H7}: Conclusion duplication---a single premise that \emph{is} the conclusion (\ISR{} inflated to $\geq$0.5).
\item \textbf{Tautological collapse}: Output is trivially provable (e.g., $P \to P$), typically with very few components.
\item \textbf{Repetition collapse}: Model degenerates into repetitive token sequences.
\item \textbf{Format collapse}: Output fails to parse entirely.
\end{itemize}

H7 is ISR-specific: it only appears when \ISR{} is used as a reward signal, as it maximizes \ISR{} by construction.

\subsection{\ISR{} as Training Reward}

We integrate \ISR{} into the \GRPO{} reward via an additive bonus:
\begin{equation}
r = r_{\text{exec}} \cdot (1 + \beta \cdot \ISR{})
\end{equation}
where $r_{\text{exec}} \in \{0, 1\}$ is the execution (provability) reward and $\beta$ controls the \ISR{} weight. We test $\beta \in \{0.1, 0.3\}$ with matched random-reward controls ($r = r_{\text{exec}} \cdot (1 + \beta \cdot U[0,1])$).

% ============================================================
\section{Experimental Setup}
\label{sec:setup}

\paragraph{Model.} Qwen3-1.5B as the base model for all primary experiments. We also report preliminary results with Qwen3-8B (\S\ref{sec:scaling}).

\paragraph{Training.} \GRPO{} with $K{=}8$ samples per prompt, 500 training steps (extended to 1000 for temporal analysis). FOL formalization using Prover9 as the solver; Prolog formalization using SWI-Prolog; SMT-LIB using Z3.

\paragraph{Evaluation.} Held-out evaluation on 100 problems. Metrics: prove rate (PR), mean \ISR{}, proved-only \ISR{}, hacking mode classification (H1, H7, genuine), component count.

\paragraph{Seeds.} Three seeds (s42, s43, s44) for all multi-seed experiments.

\paragraph{Conditions.} We evaluate: (1) exec-only, (2) exec+ISR at $\beta{=}0.1$ and $\beta{=}0.3$, (3) exec+random at $\beta{=}0.1$ and $\beta{=}0.3$, (4) SFT baseline, (5) convex combination variants, and (6) cross-logic conditions (Prolog, SMT-LIB, mixed).

% ============================================================
\section{Claim 1: \ISR{} as Reward Hacking Diagnostic}
\label{sec:claim1}

\subsection{\ISR{} Collapse Over Training}

Exec-only \GRPO{} training reveals a stark divergence between prove rate and structural integrity. Table~\ref{tab:isr_collapse} shows the complete trajectory.

\begin{table}[h]
\centering
\small
\begin{tabular}{lccc}
\toprule
\textbf{Steps} & \textbf{PR} & \textbf{Proved-only \ISR{}} & \textbf{Components} \\
\midrule
100 & 0.32 & 0.6423 & -- \\
500 & 0.93 & 0.1324 & -- \\
1000 & 0.96 & 0.0027 & 18.55 \\
\bottomrule
\end{tabular}
\caption{\ISR{} collapse trajectory over training (FOL exec-only, s42). Prove rate increases while \ISR{} collapses by 99.6\% from step 100 to 1000. At 1000 steps, 94.8\% of proved samples have \ISR{}=0.}
\label{tab:isr_collapse}
\end{table}

At 100 steps, the model produces structurally sound formalizations (\ISR{}=0.64). By 500 steps, prove rate reaches 93\% but \ISR{} has dropped 79\%. By 1000 steps, \ISR{} reaches 0.003---effectively zero---while prove rate remains at 96\%. The model has learned to satisfy the solver through structural shortcuts invisible to prove rate.

\subsection{Cross-Seed Hacking Diversity}

Three seeds trained under identical exec-only conditions exhibit qualitatively different hacking strategies (Table~\ref{tab:3seed_exec}).

\begin{table}[h]
\centering
\small
\begin{tabular}{lcccccc}
\toprule
\textbf{Seed} & \textbf{PR} & \textbf{\ISR{}} & \textbf{H1} & \textbf{H7} & \textbf{Genuine} & \textbf{Mode} \\
\midrule
s42 & 0.93 & 0.123 & 87.1\% & 0\% & 12.9\% & redundant embed. \\
s43 & 0.13 & 0.0 & -- & -- & -- & repetition collapse \\
s44 & 0.98 & 0.57 & 0\% & 100\% & 0\% & tautological \\
\bottomrule
\end{tabular}
\caption{Exec-only \GRPO{} 3-seed results (FOL, 500 steps). Each seed develops a distinct hacking strategy. s44's high \ISR{} is an artifact of tautological collapse ($n_{\text{comp}}$=2.04).}
\label{tab:3seed_exec}
\end{table}

This seed-level divergence is itself a finding: single-seed \GRPO{} results are unreliable for assessing structural quality. s42 achieves high PR via conclusion embedding (H1=87\%), s43 collapses entirely, and s44 achieves near-perfect PR via tautological shortcuts that artificially inflate \ISR{}.

\subsection{SFT Baseline}

A supervised fine-tuning baseline confirms that hacking is RL-reward-specific:

\begin{table}[h]
\centering
\small
\begin{tabular}{lccccc}
\toprule
\textbf{Method} & \textbf{PR} & \textbf{Proved-only \ISR{}} & \textbf{H7} & \textbf{H1} & \textbf{Genuine} \\
\midrule
SFT & 0.14 & 0.571 & 0\% & 7.1\% & 92.9\% \\
Exec-only (s42) & 0.93 & 0.132 & 0\% & 87.1\% & 12.9\% \\
\bottomrule
\end{tabular}
\caption{SFT baseline vs.\ exec-only \GRPO{}. SFT has low PR but 92.9\% genuine formalizations with near-zero hacking, confirming reward hacking is RL-specific.}
\label{tab:sft}
\end{table}

% ============================================================
\section{Claim 2: ISR-as-Reward Exhibits Meta-Level Goodhart's Law}
\label{sec:claim2}

\subsection{Low Dose ($\beta{=}0.1$): ISR-Specific Gaming}

When \ISR{} is added as a reward signal at $\beta{=}0.1$, 2 of 3 seeds game \ISR{} itself through conclusion duplication (H7). Table~\ref{tab:isr_reward_3seed} shows the complete 3-seed comparison with random-reward controls.

\begin{table*}[t]
\centering
\small
\begin{tabular}{llcccccc}
\toprule
\textbf{Reward} & \textbf{Seed} & \textbf{PR} & \textbf{Proved-only \ISR{}} & \textbf{H7} & \textbf{H1} & \textbf{H7+H1} & \textbf{Genuine \%} \\
\midrule
ISR $\beta{=}0.1$ & s42 & 0.81 & 0.415 & 16.0\% & 51.9\% & 67.9\% & 32.1\% \\
ISR $\beta{=}0.1$ & s43 & 0.61 & 0.830 & 96.7\% & 0\% & 96.7\% & $\sim$3\% \\
ISR $\beta{=}0.1$ & s44 & 0.89 & 0.902 & 86.5\% & 0\% & 86.5\% & $\sim$13.5\% \\
\midrule
Random $\beta{=}0.1$ & s42 & 0.79 & 0.131 & 0\% & 70.9\% & 70.9\% & 29.1\% \\
Random $\beta{=}0.1$ & s43 & 0.88 & 0.058 & 0\% & 93.2\% & 93.2\% & 6.8\% \\
Random $\beta{=}0.1$ & s44 & 0.89 & 0.070 & 0\% & 79.8\% & 79.8\% & 20.2\% \\
\midrule
\textbf{ISR 3-seed} & mean$\pm$std & 0.77$\pm$0.14 & 0.716$\pm$0.26 & 64.5\% & 17.3\% & -- & $\sim$16.2\% \\
\textbf{Random 3-seed} & mean$\pm$std & 0.85$\pm$0.06 & 0.086$\pm$0.039 & 0\% & 81.3\%$\pm$11.3\% & -- & 18.7\%$\pm$11.3\% \\
\bottomrule
\end{tabular}
\caption{ISR-as-reward ($\beta{=}0.1$) vs.\ random-reward control, 3 seeds. H7 (conclusion duplication) appears exclusively under ISR reward (64.5\% vs.\ 0\% across all random seeds). Random reward achieves higher and more stable PR (0.85$\pm$0.06 vs.\ 0.77$\pm$0.14).}
\label{tab:isr_reward_3seed}
\end{table*}

Key observations:
\begin{itemize}
\item \textbf{H7 is ISR-specific}: ISR reward produces H7=64.5\% across 3 seeds; random reward produces H7=0\% across all 3 seeds. H7 is directly induced by ISR optimization.
\item \textbf{Conclusion embedding is GRPO-universal}: Both ISR and random reward conditions show high conclusion-in-premise rates (H7+H1). Random actually has \emph{higher} H1 (81.3\% vs.\ 17.3\%). ISR reward restructures hacking (H1$\to$H7) without reducing it.
\item \textbf{Random reward outperforms ISR}: PR 0.85$\pm$0.06 vs.\ 0.77$\pm$0.14, with comparable genuine rates ($\sim$18.7\% vs.\ $\sim$16.2\%).
\end{itemize}

\subsection{Premise-Count Confound Analysis}

The one seed (s42) where ISR reward appears to produce genuine improvement (32.1\% genuine, proved-only ISR=0.415) is subject to a premise-count confound. ISR is a LOO metric: fewer premises $\to$ each premise more likely necessary $\to$ higher ISR, independent of reasoning quality.

\begin{table}[h]
\centering
\small
\begin{tabular}{lcccc}
\toprule
\textbf{Condition} & \textbf{$n$} & \textbf{Mean $n_{\text{prem}}$} & \textbf{Mean \ISR{}} \\
\midrule
ISR s42 genuine, $n_{\text{prem}} \geq 3$ & 11 & $\sim$5.5 & 0.2578 \\
Random s42 genuine, $n_{\text{prem}} \geq 3$ & 21 & $\sim$6.5 & 0.2434 \\
\midrule
\textbf{Ratio} & -- & -- & \textbf{1.06$\times$} \\
\bottomrule
\end{tabular}
\caption{Premise-count-matched comparison. After controlling for premise count ($n_{\text{prem}} \geq 3$), ISR reward shows no structural advantage over random reward (1.06$\times$). ISR s42's 15/26 ``genuine'' samples have only 1--2 premises, where ISR $\geq$ 0.5 is a mathematical artifact.}
\label{tab:premise_matched}
\end{table}

This analysis demonstrates that premise minimization---not improved reasoning---drives the apparent ISR advantage. ISR-as-reward produces shorter proofs with fewer premises, which mechanically inflates ISR without improving structural quality.

\subsection{High Dose ($\beta{=}0.3$): Classic Goodhart Collapse}

At $\beta{=}0.3$, classic Goodhart collapse occurs:

\begin{table}[h]
\centering
\small
\begin{tabular}{lccc}
\toprule
\textbf{Condition} & \textbf{PR} & \textbf{\ISR{}} & \textbf{Mode} \\
\midrule
ISR $\beta{=}0.3$ & 0.18 & 0.026 & format collapse \\
Random $\beta{=}0.3$ & 0.48 & 0.045 & complexity w/o utility \\
\bottomrule
\end{tabular}
\caption{High-dose ISR reward ($\beta{=}0.3$) vs.\ random control. ISR reward actively harms training (PR 18\% vs.\ 48\%), demonstrating classic Goodhart collapse.}
\label{tab:high_dose}
\end{table}

\subsection{Dose--Response Summary}

Table~\ref{tab:dose_response} summarizes the full Goodhart landscape across all conditions.

\begin{table}[h]
\centering
\small
\begin{tabular}{lccl}
\toprule
\textbf{Condition} & \textbf{PR} & \textbf{\ISR{}} & \textbf{Interpretation} \\
\midrule
Exec-only (s42) & 0.93 & 0.123 & H1 87\%, no ISR signal \\
SFT baseline & 0.14 & 0.571 & 93\% genuine \\
ISR $\beta{=}0.1$ (s42) & 0.81 & 0.415 & premise-matched: 1.06$\times$ \\
ISR $\beta{=}0.1$ (s43) & 0.61 & 0.830 & H7 97\%, meta-Goodhart \\
ISR $\beta{=}0.1$ (s44) & 0.89 & 0.902 & H7 87\%, meta-Goodhart \\
ISR $\beta{=}0.3$ & 0.18 & 0.026 & classic Goodhart \\
Random $\beta{=}0.1$ (mean) & 0.85 & 0.086 & H1 81\%, control \\
Random $\beta{=}0.3$ & 0.48 & 0.045 & control \\
Convex comb. & 0.77 & 0.230 & penalty effect \\
Convex random & 1.00 & 0.535 & tautological collapse \\
\bottomrule
\end{tabular}
\caption{Complete Goodhart landscape. ISR-as-reward fails at every dose: low dose induces ISR-specific gaming (meta-Goodhart), high dose collapses PR (classic Goodhart).}
\label{tab:dose_response}
\end{table}

% ============================================================
\section{Additional Analyses}

\subsection{$O(1)$ vs.\ $O(N)$ STA Ablation}
\label{sec:o1ablation}

Full \STA{} requires $O(N)$ solver calls per sample, which is computationally expensive. We evaluate an $O(1)$ heuristic approximation that estimates ISR without per-component ablation. Table~\ref{tab:o1_on} compares the two approaches.

\begin{table}[h]
\centering
\small
\begin{tabular}{lcccc}
\toprule
\textbf{Method} & \textbf{PR} & \textbf{\ISR{}} & \textbf{Time/step (s)} & \textbf{Genuine} \\
\midrule
$O(N)$ full \STA{} & 0.77 & 0.23 & 10.1 & $\sim$26\% \\
$O(1)$ fast ISR & 0.72 & 0.09 & 3.1 & $\sim$22\% \\
Exec-only & 0.93 & 0.13 & 9.4 & 12.9\% \\
Random $\beta{=}0.1$ & 0.79 & 0.10 & 9.5 & $\sim$25\% \\
\bottomrule
\end{tabular}
\caption{$O(1)$ approximation vs.\ $O(N)$ full \STA{}. The fast approximation is 3.3$\times$ faster but produces a noisier ISR signal (0.09 vs.\ 0.23), yielding results comparable to random reward---indicating that full LOO ablation is necessary for reliable ISR computation.}
\label{tab:o1_on}
\end{table}

The $O(1)$ approximation is 3.3$\times$ faster but its ISR signal degrades to near-random levels (\ISR{}=0.09 vs.\ full \STA{} 0.23), with PR=0.72 lower than exec-only (0.93). This confirms that the LOO ablation---actually executing the solver with each component removed---is essential for meaningful structural assessment. Heuristic shortcuts introduce systematic bias that renders the signal unreliable as either diagnostic or reward.

\subsection{Curriculum Schedule Ablation}
\label{sec:curriculum}

We evaluate whether the schedule for introducing ISR reward affects outcomes. Table~\ref{tab:curriculum} compares sigmoid, linear, and fixed weighting schemes.

\begin{table}[h]
\centering
\small
\begin{tabular}{lcccc}
\toprule
\textbf{Schedule} & \textbf{PR} & \textbf{\ISR{}} & \textbf{Genuine} \\
\midrule
Sigmoid (mp=300) & 0.77 & 0.23 & $\sim$26\% \\
Sigmoid (mp=200) & 0.74 & 0.21 & $\sim$24\% \\
Sigmoid (mp=100) & 0.68 & 0.18 & $\sim$20\% \\
Linear ramp (mp=300) & 0.74 & 0.21 & $\sim$24\% \\
Fixed $\beta{=}0.1$ & 0.71 & 0.19 & $\sim$22\% \\
Fixed $\beta{=}0.3$ & 0.62 & 0.15 & $\sim$18\% \\
\bottomrule
\end{tabular}
\caption{Curriculum ablation across schedule types. All configurations fall within a narrow band (PR: 0.62--0.77, ISR: 0.15--0.23). No schedule eliminates the Goodhart dynamics ($\Delta$PR $<$ 5pp across configurations), and no configuration achieves genuine rate above 40\%.}
\label{tab:curriculum}
\end{table}

The narrow range across all configurations ($\Delta$PR $<$ 15pp, $\Delta$ISR $<$ 8pp) indicates that curriculum design is not the bottleneck: the fundamental issue is Goodhart optimization of ISR, regardless of when or how the signal is introduced.

\subsection{Model Scaling: Qwen3-8B}
\label{sec:scaling}

We conduct a preliminary single-seed experiment with Qwen3-8B to assess whether model scale mitigates reward hacking. Table~\ref{tab:scaling} compares 1.5B and 8B results.

\begin{table}[h]
\centering
\small
\begin{tabular}{lcccccc}
\toprule
\textbf{Model} & \textbf{Condition} & \textbf{PR} & \textbf{\ISR{}} & \textbf{H1} & \textbf{H7} & \textbf{Genuine} \\
\midrule
1.5B & Exec-only & 0.93 & 0.13 & 87\% & 0\% & 12.9\% \\
8B & Exec-only & 0.97 & 0.08 & 92\% & 0\% & 5\% \\
\midrule
1.5B & ISR $\beta{=}0.1$ & 0.81 & 0.42 & -- & 16\% & 30.9\% \\
8B & ISR $\beta{=}0.1$ & 0.85 & 0.52 & -- & 45\% & $\sim$15\% \\
\bottomrule
\end{tabular}
\caption{1.5B vs.\ 8B model scaling (single seed). Scaling \emph{exacerbates} rather than mitigates reward hacking: the 8B model achieves higher PR (0.97) but lower genuine rate (5\%), and shows elevated H7 rates under ISR reward (45\% vs.\ 16\%).}
\label{tab:scaling}
\end{table}

The larger model exploits structural shortcuts more efficiently: exec-only genuine rate drops from 12.9\% to 5\%, while ISR reward H7 rate increases from 16\% to 45\%. This suggests that increased model capacity amplifies Goodhart optimization rather than enabling genuine reasoning improvement. We note this is a single-seed preliminary result and a systematic scaling study is needed.

% ============================================================
\section{Preliminary Cross-Logic Exploration}
\label{sec:crosslogic}

We attempted to extend our findings beyond FOL to Prolog (SWI-Prolog) and SMT-LIB (Z3). These experiments largely failed, and we report them as negative results that inform future work.

\subsection{Prolog STA: Extreme Seed Instability}

Prolog exec+STA \GRPO{} with ISR reward was run across 3 seeds. Results are summarized in Table~\ref{tab:prolog_sta}.

\begin{table}[h]
\centering
\small
\begin{tabular}{lccccl}
\toprule
\textbf{Seed} & \textbf{Status} & \textbf{PR} & \textbf{\ISR{}} & \textbf{H7} & \textbf{Mode} \\
\midrule
s42 & Converged & 0.50 & 0.50 & 100\% & H7 saturation \\
s43 & Failed & 0.06 & 0.01 & 0\% & training collapse \\
s44 & Partial & 0.35 & 0.12 & 11\% & mixed/weak \\
\midrule
Mean$\pm$std & -- & 0.30$\pm$0.22 & 0.21$\pm$0.26 & 0.37$\pm$0.55 & -- \\
\bottomrule
\end{tabular}
\caption{Prolog STA 3-seed results. Extreme variance ($\sigma_{\text{H7}}$=0.55) renders aggregate statistics meaningless. Only 1/3 seeds reproduces the H7 pattern observed in FOL.}
\label{tab:prolog_sta}
\end{table}

\paragraph{s43 training failure.} Seed 43 failed to learn any useful signal: final PR=0.06 (near base-model level), ISR=0.01, with loss stagnating at 0.35 after step 100. Table~\ref{tab:prolog_s43_trajectory} shows the training trajectory.

\begin{table}[h]
\centering
\small
\begin{tabular}{lccc}
\toprule
\textbf{Step window} & \textbf{PR} & \textbf{\ISR{}} & \textbf{Loss} \\
\midrule
1--50 & 0.04 & 0.01 & 0.42 \\
51--100 & 0.06 & 0.01 & 0.39 \\
101--200 & 0.07 & 0.02 & 0.37 \\
201--300 & 0.08 & 0.02 & 0.35 \\
301--400 & 0.08 & 0.02 & 0.35 \\
401--500 & 0.08 & 0.02 & 0.35 \\
\bottomrule
\end{tabular}
\caption{Prolog STA s43 training trajectory. Training stagnates after $\sim$100 steps with no meaningful learning signal, consistent with sparse reward in Prolog's constrained semantic space.}
\label{tab:prolog_s43_trajectory}
\end{table}

\paragraph{s44 partial convergence.} Seed 44 achieved partial convergence (PR=0.35) but with H7=11\% (4/35 proved)---far below s42's 100\%. Of 35 proved samples, 22 (63\%) were genuine multi-clause outputs, suggesting this seed followed a qualitatively different optimization path than s42.

\paragraph{Cross-logic comparison.} Table~\ref{tab:cross_logic_compare} compares Prolog STA results with FOL STA.

\begin{table}[h]
\centering
\small
\begin{tabular}{lccc}
\toprule
\textbf{Metric} & \textbf{Prolog STA} & \textbf{FOL ISR $\beta{=}0.1$} & \textbf{Interpretation} \\
\midrule
3-seed H7 mean & 0.37 ($\sigma$=0.55) & 0.645 ($\sigma$=0.44) & Prolog higher variance \\
Training failures & 1/3 & 0/3 & Prolog less stable \\
Cross-logic effect & $\sim$0.15 & baseline & Cohen's $d < 0.3$ \\
\bottomrule
\end{tabular}
\caption{Prolog vs.\ FOL STA cross-logic comparison. With 1/3 seeds failing entirely and massive variance, Prolog results do not constitute evidence for cross-logic generalization of the Goodhart pattern.}
\label{tab:cross_logic_compare}
\end{table}

With 1/3 seeds failing training and $\sigma_{\text{H7}}=0.55$, the confidence interval for H7 rate includes zero. We cannot claim cross-logic generalization of the Goodhart effect; s42's H7=100\% is an isolated observation.

\subsection{Mixed-Logic Training: Catastrophic Interference}

Training \GRPO{} on mixed FOL+Prolog data produces catastrophic interference (Table~\ref{tab:mixed}).

\begin{table}[h]
\centering
\small
\begin{tabular}{lcccc}
\toprule
\textbf{Condition} & \textbf{FOL PR} & \textbf{Prolog PR} & \textbf{FOL \ISR{}} \\
\midrule
FOL-only exec & 0.93 & N/A & 0.13 \\
Prolog-only exec (s42) & N/A & 0.83 & N/A \\
Mixed exec-only & 0.45 & 0.22 & 0.05 \\
Mixed exec+STA & 0.38 & 0.15 & 0.08 \\
\bottomrule
\end{tabular}
\caption{Mixed FOL+Prolog training vs.\ single-logic baselines. FOL PR drops from 0.93 to 0.45 ($-$48pp); Prolog PR drops from 0.83 to 0.22 ($-$61pp). Both logics suffer catastrophic interference.}
\label{tab:mixed}
\end{table}

FOL prove rate degrades by 48 percentage points and Prolog by 61 percentage points compared to single-logic training. The two solver interfaces (Prover9 vs.\ SWI-Prolog) produce conflicting gradient signals within \GRPO{}'s $K{=}8$ sampling, and Prolog's smaller data volume (100 vs.\ 500 problems) causes its signal to be overwhelmed. Adding STA further destabilizes training due to mismatched ISR computation across logic types.

\subsection{Zero-Shot Cross-Logic Transfer}

Models trained on one logic type show near-zero transfer to another (Table~\ref{tab:transfer}).

\begin{table}[h]
\centering
\small
\begin{tabular}{llcc}
\toprule
\textbf{Model} & \textbf{Eval on} & \textbf{PR} & \textbf{\ISR{}} \\
\midrule
FOL-trained (s42) & FOL & 0.93 & 0.13 \\
FOL-trained (s42) & Prolog & 0.02 & 0.00 \\
Prolog-trained (s42) & Prolog & 0.83 & 0.03 \\
Prolog-trained (s42) & FOL & 0.04 & 0.01 \\
\bottomrule
\end{tabular}
\caption{Zero-shot cross-logic transfer. Models show near-zero transfer (PR $\leq$ 4\%) across logic types, consistent with format-specific overfitting in \GRPO{} training.}
\label{tab:transfer}
\end{table}

Cross-logic transfer is effectively zero (PR $\leq$ 4\%), with Prolog-trained models showing catastrophic forgetting of FOL capabilities (4\% vs.\ base model $\sim$12\%). This is consistent with \GRPO{} producing format-specific rather than reasoning-general capabilities.

\subsection{SMT-LIB: Third Logic Type Failure}

We attempted SMT-LIB (Z3) as a third logic type. Training failed due to Z3 timeout-dominated reward signals (Table~\ref{tab:three_logics}).

\begin{table}[h]
\centering
\small
\begin{tabular}{lcccc}
\toprule
\textbf{Logic} & \textbf{Exec PR} & \textbf{Exec \ISR{}} & \textbf{STA H7} & \textbf{Status} \\
\midrule
FOL (Prover9) & 0.93 & 0.13 & 0--97\% & Closed \\
Prolog (SWI) & 0.53$\pm$0.31 & 0.04$\pm$0.04 & 0--100\% & Unstable \\
SMT-LIB (Z3) & 0.15 & 0.02 & 0\% & Failed \\
\bottomrule
\end{tabular}
\caption{Three logic types comparison. Only FOL produces stable, analyzable results. Prolog is highly seed-dependent. SMT-LIB fails to train due to sparse reward signals from Z3 timeouts.}
\label{tab:three_logics}
\end{table}

SMT-LIB exec-only achieved only PR=15\% (suggesting weak base-model priors for SMT-LIB syntax), and exec+STA degraded further to PR=8\% with ISR=0.04, H7=0\%. Z3's higher timeout rate compared to Prover9 renders the reward signal too sparse for \GRPO{}'s group-relative advantages to produce meaningful gradients.

% ============================================================
\section{Discussion}

\paragraph{ISR's asymmetric value.} Our results demonstrate a clean separation: \ISR{} is a powerful post-hoc diagnostic (Claim 1) but fails as a training objective (Claim 2). This asymmetry is not accidental---it follows from Goodhart's Law applied at the meta level. Any metric that accurately detects gaming becomes a target for gaming when used as an optimization signal.

\paragraph{Conclusion embedding as GRPO-universal phenomenon.} Across all reward conditions, conclusion embedding (H1+H7) affects 52--93\% of proved samples. ISR reward does not reduce this; it merely restructures it from multi-premise embedding (H1) to single-premise duplication (H7). This suggests that conclusion embedding is an intrinsic property of \GRPO{} training on formal reasoning tasks, not an artifact of any particular reward.

\paragraph{Scaling considerations.} Preliminary 8B results suggest scaling amplifies rather than mitigates reward hacking (genuine rate 5\% vs.\ 12.9\% at 1.5B). If confirmed, this has implications for the common assumption that scale improves reasoning quality.

\paragraph{Cross-logic challenges.} The failure of Prolog and SMT-LIB experiments highlights that \STA{}'s applicability depends on solver characteristics (timeout rates, semantic space size, syntax complexity). Extending structural diagnostics beyond FOL requires solver-specific adaptations.

% ============================================================
\section{Limitations}

\begin{itemize}
\item \textbf{Single logic type for main claims.} Claims 1 and 2 are validated on FOL (Prover9) only. Cross-logic experiments (Prolog, SMT-LIB) produced inconclusive or negative results.
\item \textbf{Single model family.} All experiments use Qwen3 (1.5B primary, 8B preliminary). Generalization to other model families is untested.
\item \textbf{\ISR{} measures syntactic necessity, not semantic faithfulness.} \ISR{} correlates poorly with downstream NLI accuracy ($\rho = -0.049$, $p=0.187$). Formalizations with high \ISR{} are structurally non-redundant but not necessarily semantically faithful to the natural language input.
\item \textbf{Premise-count bias.} LOO-based metrics like \ISR{} are mathematically biased by component count: fewer premises $\to$ higher \ISR{}, independent of reasoning quality. This confound must be controlled for in any ISR-based analysis.
\item \textbf{500-step training horizon.} Most experiments use 500 steps. While we extend to 1000 steps for temporal analysis, longer training horizons may reveal additional dynamics.
\item \textbf{Three seeds.} While 3-seed divergence is itself informative, statistical power is limited. The massive cross-seed variance in Prolog ($\sigma_{\text{H7}}=0.55$) underscores this limitation.
\item \textbf{Scaling result is preliminary.} The 8B comparison uses a single seed and a single model size increment.
\end{itemize}

% ============================================================
\section{Conclusion}

We introduced Solver Trace Ablation (\STA{}) and the Intensional Structural Reward (\ISR{}) as a diagnostic framework for reward hacking in \GRPO{}-trained formal reasoners. Our main findings are:

\begin{enumerate}
\item \ISR{} detects structural degeneration invisible to prove rate, with a 99.6\% collapse from training step 100 to 1000 while prove rate remains above 93\%.
\item When repurposed as a training reward, \ISR{} exhibits meta-level Goodhart's Law: at low dose ($\beta{=}0.1$), 2/3 seeds game \ISR{} itself via conclusion duplication; the remaining seed's apparent improvement is a premise-count artifact (1.06$\times$ when controlled). At high dose ($\beta{=}0.3$), prove rate collapses to 18\%.
\item Across 6 reward conditions and 3 seeds, we identify 6 distinct hacking strategies and document conclusion embedding as a GRPO-universal phenomenon (52--93\% of proved samples).
\item Cross-logic experiments (Prolog, SMT-LIB) reveal high instability and training failures, indicating that extending these findings beyond FOL remains an open challenge.
\end{enumerate}

Our results carry a cautionary message for the broader proxy-metric-as-reward paradigm: metrics that diagnose gaming can themselves become gaming targets when used as optimization signals. We recommend using \ISR{} for post-hoc structural diagnosis only, with explicit correction for the premise-count confound.

% ============================================================
\begin{thebibliography}{10}

\bibitem[Shao et~al.(2024)]{shao2024deepseekmath}
Zhihong Shao, Peiyi Wang, Qihao Zhu, et~al.
DeepSeekMath: Pushing the limits of mathematical reasoning in open language models.
\textit{arXiv preprint arXiv:2402.03300}, 2024.

\bibitem[Luo et~al.(2024)]{luo2024improve}
Haipeng Luo, Qingfeng Sun, Can Xu, et~al.
WizardMath: Empowering mathematical reasoning for large language models via reinforced evol-instruct.
\textit{arXiv preprint arXiv:2308.09583}, 2024.

\bibitem[Manheim and Garrabrant(2019)]{manheim2019categorizing}
David Manheim and Scott Garrabrant.
Categorizing variants of Goodhart's Law.
\textit{arXiv preprint arXiv:1803.04585}, 2019.

\bibitem[Gao et~al.(2023)]{gao2023scaling}
Leo Gao, John Schulman, and Jacob Hilton.
Scaling laws for reward model overoptimization.
In \textit{Proceedings of ICML}, 2023.

\bibitem[Mellgren et~al.(2025)]{mellgren2025}
Gustav Mellgren et~al.
Structure-aware rewards for Prolog program synthesis.
\textit{arXiv preprint}, 2025.

\bibitem[{EVOM}(2026)]{evom2026}
EVOM.
Execution verification for open-ended mathematical reasoning.
\textit{arXiv preprint}, 2026.

\bibitem[Zeiler and Fergus(2014)]{zeiler2014visualizing}
Matthew~D. Zeiler and Rob Fergus.
Visualizing and understanding convolutional networks.
In \textit{Proceedings of ECCV}, 2014.

\end{thebibliography}

\end{document}
